\paragraph*{哪些区间需要建点？}
给定排列 $a$，下标从 $0$ 开始。若 $a_l,\ldots,a_r$ 的值恰好填满一段连续整数，
称 $[l,r]$ 为\textbf{连续段}。由于元素互异，判据是 $\max-\min=r-l$。
两个区间相交且互不包含时称为交叉。与所有连续段都不交叉的连续段称为
\textbf{本原段}（强区间）。把每个本原段挂到严格包含它的最小本原段下面，
就得到析合树。单点是叶子，整个排列是根，儿子按原下标从左到右排列。

把每个儿子缩成一个块，按其值域从小到大编号。若各块编号从左到右递增或递减，
父亲就是\textbf{合点}，任意一段相邻儿子都能拼成连续段。
例如 $(2,1)\mid(4,3)\mid(6,5)$ 的三块值域依次为 $[1,2],[3,4],[5,6]$，
根是上升合点，前两块就能拼成一个连续段。另一类父亲是\textbf{析点}，
跨越多个儿子的连续段只能取全部儿子，例如 $(2,4,1,3)$。
任一连续段若跨过本原段的边界，就必须把该本原段取全。
因此，树上除结点本身外的连续段，都可以在某个合点处用连续的完整儿子表示。

\paragraph*{固定右端点，找最短的连续后缀}
扫到下标 $i$ 时，线段树的每个叶子代表一个候选左端点 $j\le i$，维护值域里缺了多少个数：
\[
  F_i(j)=\max_{j\le h\le i} a_h-\min_{j\le h\le i} a_h-(i-j).
\]
$F_i(j)\ge 0$，等于零就说明 $[j,i]$ 连续。若当前块 $p$ 已覆盖 $[l_p,i]$，
要向左扩大它，调用 \texttt{findlast}$(l_p-1)$，在 $[0,l_p-1]$ 中找最右的零。
线段树维护区间最小值，最小值大于零就跳过；下传标记后先搜右儿子。
同一个 $i$ 下可以合并多次，只有查询上界随 $p$ 改变。

从 $i-1$ 扫到 $i$，先给所有旧起点 $[0,i-1]$ 减一。以最大值栈为例，
若 $a_i>a_x$ 使栈顶 $x$ 弹出，记弹出后的栈顶为 $b$（空栈取 $-1$）。
旧最大值 $a_x$ 恰好管左端点 $[b+1,x]$：再往左会包含更大的 $a_b$，
再往右就不含 $a_x$。代码在这一块减去 $a_x$，并用 $R$ 记录已撤掉的最左位置。
这些块首尾相接，弹完后在 $[R,i]$ 统一加上 $a_i$。最小值栈对称地先加回旧最小值，
再减去新最小值；新起点的两项正好抵消，$F_i(i)=0$。

\paragraph*{森林栈怎样合并？}
栈 \texttt{stk} 从左到右存放已处理前缀的极大连续块。令 $p$ 为新叶子 $i$，
$q$ 为紧邻其左的栈顶，反复尝试以下两步。

\begin{enumerate}
  \item \textbf{先追加。} 若 $q$ 为上升合点且 $\max_q+1=\min_p$，
        或为下降合点且 $\min_q-1=\max_p$，就弹出 $q$，把 $p$ 接为其末子，
        令 $p=q$ 继续。例如 $(1,2)$ 后接 $3$ 时，$(2,3)$ 与旧块 $(1,2)$ 交叉，
        所以三个叶子应同属一个上升合点。
  \item \textbf{再建父亲。} 追加失败后查最右零点 $L$。找不到就把 $p$ 压栈并结束本轮。
        否则 $L$ 落在某个完整栈块的左端。找到满足 $l_{\texttt{stk}[k]}=L$ 的栈下标 $k$，
        依次取出从 $k$ 到栈顶的各块，最后接上 $p$，把新父亲作为 $p$ 继续合并。
\end{enumerate}

\textbf{新父亲的类型。} 若 $L=l_q$，只有 $q,p$ 两个儿子，其值域紧邻，
比较两块的最小值就能定升降。若 $L<l_q$，至少取了两棵旧树再加 $p$，
新父亲必为析点。因为只取旧块时，若相邻几块能合并，早已不会分留在栈中；
若取 $p$ 和部分旧块仍连续，又会得到比 $L$ 更靠右的零。析点至少有四个儿子：
三个块里，值域排名居中的块总能和一个邻块拼成连续段。

\textbf{手推一轮。} 对 $(2,1,4,6,3,5)$ 插入最后的 $5$，此时栈为
$(2,1)\mid(4)\mid(6)\mid(3)$。向左试，$(3,5)$ 和 $(6,3,5)$ 都缺 $4$，
而 $(4,6,3,5)$ 的值域是 $[3,6]$。故 $L=2$，对应栈下标 $k=1$，
后三棵旧树和 $5$ 先组成一个析点。仍在这一轮，把这个析点当作 $p$，
再与 $(2,1)$ 合并，两块值域依次为 $[1,2],[3,6]$，新根是上升合点。

\paragraph*{接口与应用}
每组用新对象调用 \texttt{DivideCombineTree tr; tr.build(a);}。
非空排列的叶子编号等于下标；在 \texttt{tr.t} 中找 \texttt{fa == -1} 即得根。
\texttt{l,r} 是下标闭区间，\texttt{mi,ma} 是值域两端，\texttt{son} 按下标有序。
\texttt{type} 取 $-1,0,1,2$，依次为叶子、析点、上升合点、下降合点。
两个单调栈各至多弹出 $n$ 次，森林中取出的栈后缀随即移除，
故建树耗时 $O(n\log n)$，至多 $2n-1$ 个结点，总空间 $O(n)$。

\textbf{统计连续段。} 按包含它的最小树结点计数：每个叶子或析点贡献 $1$，
有 $d$ 个儿子的合点贡献 $\binom d2$，即选首尾两个不同儿子。总数用 64 位整数。

\textbf{求包含 $[l,r]$ 的最小连续段。} 单点直接返回；否则设 $u$ 是两端叶子的 LCA。
若 $u$ 是析点，取整个 $u$；若是合点，找到 $l,r$ 所在的两个儿子，
取左儿子的左端点到右儿子的右端点。
